% Hand-drawn circuitikz schematic of the control amplifier core.
% NOT generated by maestro2tex -- the tool only places it, via the "order" list in
% configs/anatop_tb_controlamp_cell.json. The `sch_` prefix keeps it clear of the
% generated fig_*/tab_* namespace.
%
% Devices and connectivity transcribed from subckt OpAmpCurrentMirrorP in
% simarchive/anatop_tb_controlamp/cm_2026-07-27/Plan.4.Run.2/psf/
%   UnityFeedbackWBiasTB/netlist/input.scs   (Spectre order M(drain gate source bulk)).
%
% Placement and device orientation follow the Virtuoso schematic itself, read from
% castalia/OpAmpCurrentMirrorP/schematic via dbOpenCellViewByType, so this figure
% reads the same way as the cellview:
%   column x=0.75  M1L, M2L    column x=3.25  M1R, M2R   (input pair + unit diodes)
%   column x=4.75  M4L, M2LM   column x=6.75  M4R, M2RM  (the two output branches)
%   rows y=1.1875 M6/M4x | 0.1875 M5x | -0.8125 M1x | -1.8125 M3x | -2.8125 M2x
% Orientation R0 -> gate on the left, MY -> mirrored, gate on the right. The input
% devices are R0 and MY respectively, so their gates face outwards as in the cellview.
%
% Matching dummies (M1dum, M4dum, M5dum, M2*dum, M3*dum) carry no signal and are not
% drawn. Identically labelled nets are connected; wires crossing without a filled dot
% are not. If the cell is re-drawn, re-read the cellview before trusting this figure.

\begin{figure}[htbp]
  \centering
  \resizebox{\linewidth}{!}{%
\begin{circuitikz}[scale=1.0, transform shape]
\ctikzset{bipoles/length=0.9cm, transistors/scale=0.78}
\tikzset{
  wire/.style={line width=0.6pt},
  nd/.style={circle, fill=black, inner sep=1.1pt},
  nlab/.style={font=\footnotesize, fill=white, inner sep=1pt},
  dlab/.style={font=\scriptsize, align=center, inner sep=1pt},
  ilab/.style={font=\scriptsize, align=center, inner sep=2pt, text=black!60},
  grp/.style={draw=black!30, dashed, line width=0.5pt, rounded corners=3pt},
}
% ---- geometry: x from the cellview columns, y from its rows -----------------
\def\xL{0}        % M1L / M2L      (cellview x=0.75)
\def\xT{2.3}      % M6 tail        (cellview x=1.75)
\def\xR{6.0}      % M1R / M2R      (cellview x=3.25)
\def\xB{9.8}      % B branch       (cellview x=4.75)
\def\xO{14.2}     % output branch  (cellview x=6.75)
\def\yvss{0.4} \def\yvdd{12.2}
\def\yNm{1.9}     % M2x  row
\def\yNc{4.1}     % M3x  row
\def\yIn{6.4}     % M1x  row
\def\yPc{8.7}     % M5x  row
\def\yPm{10.8}    % M6 / M4x row
\def\ymid{5.0}    % VG2L / VG2R gate-tap level
\def\ycpR{5.55}   % CpR corridor (mid-span: must not touch a wire endpoint)
\def\ycpL{4.65}   % CpL corridor
\def\yA{7.5}      % tail node A

% ---- rails ----------------------------------------------------------------
\draw[wire] (-2.4,\yvdd) -- (16.6,\yvdd) node[right, font=\footnotesize] {\texttt{vdd!}};
\draw[wire] (-2.4,\yvss) -- (16.6,\yvss) node[right, font=\footnotesize] {\texttt{vss!}};

% ---- group outlines --------------------------------------------------------
\node[ilab, anchor=south] at (\xT+0.4,\yvss-1.25) {input pair and unit diode loads};
\node[ilab, anchor=south] at (\xB+2.2,\yvss-1.25) {output branches: node \texttt{B}, then \texttt{VOut}};
\draw[grp] (-1.9,\yvss-0.45) -- (7.9,\yvss-0.45);
\draw[grp] (8.5,\yvss-0.45) -- (16.1,\yvss-0.45);

% =========================== tail M6 (R0) ==================================
\draw (\xT,\yPm) node[pmos](M6){};
\draw[wire] (M6.source) -- (\xT,\yvdd);
\draw[wire] (M6.drain)  -- (\xT,\yA);
\draw[wire] (M6.gate) -- ++(-0.8,0) node[nlab, anchor=east] {\texttt{bp}};
\node[dlab, anchor=west] at (\xT+0.5,\yPm) {M6\\\scriptsize $m$=1, $nf$=4};
\node[ilab, anchor=west] at (\xT+0.25,\yA+0.45) {$I_{\mathrm{tail}}$};

% =========================== input pair ====================================
\draw[wire] (\xL,\yA) -- (\xR,\yA);
\node[nd] at (\xT,\yA) {};
\node[nlab, anchor=south] at (\xT-1.0,\yA+0.06) {\texttt{A}};

\draw (\xL,\yIn) node[pmos](M1L){};                       % R0  -> gate faces left
\draw[wire] (M1L.source) -- (\xL,\yA);
\draw[wire] (M1L.drain)  -- (\xL,\yNc+0.55);
\draw[wire] (M1L.gate) -- ++(-1.0,0) node[nlab, anchor=east] {\texttt{VInP} $(+)$};
\node[dlab, anchor=west] at (\xL+0.5,\yIn) {M1L\\\scriptsize $m$=4};

\draw (\xR,\yIn) node[pmos, xscale=-1](M1R){};            % MY  -> gate faces right
\draw[wire] (M1R.source) -- (\xR,\yA);
\draw[wire] (M1R.drain)  -- (\xR,\yNc+0.55);
\draw[wire] (M1R.gate) -- ++(1.0,0) node[nlab, anchor=west] {\texttt{VInM} $(-)$};
\node[dlab, anchor=east] at (\xR-0.5,\yIn) {M1R\\\scriptsize $m$=4};

\node[nd] at (\xL,\ymid) {};
\node[nlab, anchor=east] at (\xL-0.2,\ymid) {\texttt{VG2L}};
\node[nd] at (\xR,\ymid) {};
\node[nlab, anchor=west] at (\xR+0.2,\ymid) {\texttt{VG2R}};

% =========================== unit diode loads (m=1) ========================
\draw (\xL,\yNc) node[nmos, xscale=-1](M3L){};            % MY
\draw[wire] (M3L.drain)  -- (\xL,\yNc+0.55);
\draw[wire] (M3L.source) -- (\xL,\yNm-0.55);
\draw[wire] (M3L.gate) -- ++(0.8,0) node[nlab, anchor=west] {\texttt{bnc}};
\node[dlab, anchor=east] at (\xL-0.55,\yNc) {M3L\\\scriptsize $m$=1};
\draw (\xL,\yNm) node[nmos](M2L){};                        % R0
\draw[wire] (M2L.drain)  -- (\xL,\yNm-0.55);
\draw[wire] (M2L.source) -- (\xL,\yvss);
\node[dlab, anchor=west] at (\xL+0.5,\yNm) {M2L\\\scriptsize $m$=1};
\draw[wire] (M2L.gate) -- (\xL-1.4,\yNm) -- (\xL-1.4,\ymid) -- (\xL,\ymid);

\draw (\xR,\yNc) node[nmos](M3R){};                        % R0
\draw[wire] (M3R.drain)  -- (\xR,\yNc+0.55);
\draw[wire] (M3R.source) -- (\xR,\yNm-0.55);
\draw[wire] (M3R.gate) -- ++(-0.8,0) node[nlab, anchor=east] {\texttt{bnc}};
\node[dlab, anchor=west] at (\xR+0.55,\yNc) {M3R\\\scriptsize $m$=1};
\draw (\xR,\yNm) node[nmos, xscale=-1](M2R){};             % MY
\draw[wire] (M2R.drain)  -- (\xR,\yNm-0.55);
\draw[wire] (M2R.source) -- (\xR,\yvss);
\node[dlab, anchor=east] at (\xR-0.5,\yNm) {M2R\\\scriptsize $m$=1};
\draw[wire] (M2R.gate) -- (\xR+1.4,\yNm) -- (\xR+1.4,\ymid) -- (\xR,\ymid);

% =========================== branch 1: node B ==============================
\draw (\xB,\yPm) node[pmos](M4L){};                        % R0
\draw[wire] (M4L.source) -- (\xB,\yvdd);
\draw[wire] (M4L.drain)  -- (\xB,\yPc+0.55);
\node[dlab, anchor=west] at (\xB+0.5,\yPm) {M4L\\\scriptsize $m$=4};
\draw (\xB,\yPc) node[pmos, xscale=-1](M5L){};             % MY
\draw[wire] (M5L.source) -- (\xB,\yPc+0.55);
\draw[wire] (M5L.drain)  -- (\xB,\yNc+0.55);
\draw[wire] (M5L.gate) -- ++(0.8,0) node[nlab, anchor=west] {\texttt{bpc}};
\node[dlab, anchor=east] at (\xB-0.55,\yPc) {M5L\\\scriptsize $m$=4};
% M4L is the diode of the PMOS mirror: its gate returns to B
\draw[wire] (M4L.gate) -- (\xB-1.4,\yPm) -- (\xB-1.4,\yPc-1.15) -- (\xB,\yPc-1.15);
\node[nd] at (\xB,\yPc-1.15) {};

\draw (\xB,\yNc) node[nmos, xscale=-1](M3LM){};            % MY
\draw[wire] (M3LM.drain)  -- (\xB,\yNc+0.55);
\draw[wire] (M3LM.source) -- (\xB,\yNm-0.55);
\draw[wire] (M3LM.gate) -- ++(0.8,0) node[nlab, anchor=west] {\texttt{bnc}};
\node[dlab, anchor=east] at (\xB-0.55,\yNc) {M3LM\\\scriptsize $m$=4};
\draw (\xB,\yNm) node[nmos](M2LM){};                       % R0
\draw[wire] (M2LM.drain)  -- (\xB,\yNm-0.55);
\draw[wire] (M2LM.source) -- (\xB,\yvss);
\node[dlab, anchor=west] at (\xB+0.5,\yNm) {M2LM\\\scriptsize $m$=4};
\draw[wire] (M2LM.gate) -- ++(-0.9,0) node[nlab, anchor=east] {\texttt{VG2R}};

\node[nlab, anchor=east] at (\xB-0.2,\ymid) {\texttt{B}};

% =========================== branch 2: output ==============================
\draw (\xO,\yPm) node[pmos](M4R){};                        % R0
\draw[wire] (M4R.source) -- (\xO,\yvdd);
\draw[wire] (M4R.drain)  -- (\xO,\yPc+0.55);
\draw[wire] (M4R.gate) -- ++(-0.9,0) node[nlab, anchor=east] {\texttt{B}};
\node[dlab, anchor=west] at (\xO+0.5,\yPm) {M4R\\\scriptsize $m$=4};
\draw (\xO,\yPc) node[pmos](M5R){};                        % R0
\draw[wire] (M5R.source) -- (\xO,\yPc+0.55);
\draw[wire] (M5R.drain)  -- (\xO,\yNc+0.55);
\draw[wire] (M5R.gate) -- ++(-0.8,0) node[nlab, anchor=east] {\texttt{bpc}};
\node[dlab, anchor=west] at (\xO+0.5,\yPc) {M5R\\\scriptsize $m$=4};

\draw (\xO,\yNc) node[nmos](M3RM){};                       % R0
\draw[wire] (M3RM.drain)  -- (\xO,\yNc+0.55);
\draw[wire] (M3RM.source) -- (\xO,\yNm-0.55);
\draw[wire] (M3RM.gate) -- ++(-0.8,0) node[nlab, anchor=east] {\texttt{bnc}};
\node[dlab, anchor=west] at (\xO+0.5,\yNc) {M3RM\\\scriptsize $m$=4};
\draw (\xO,\yNm) node[nmos, xscale=-1](M2RM){};            % MY
\draw[wire] (M2RM.drain)  -- (\xO,\yNm-0.55);
\draw[wire] (M2RM.source) -- (\xO,\yvss);
\node[dlab, anchor=east] at (\xO-0.5,\yNm) {M2RM\\\scriptsize $m$=4};
\draw[wire] (M2RM.gate) -- ++(0.9,0) node[nlab, anchor=west] {\texttt{VG2L}};

\node[nd] at (\xO,\ymid) {};
\draw[wire] (\xO,\ymid) -- (16.3,\ymid);
\node[nlab, anchor=south] at (15.6,\ymid+0.1) {\texttt{VOut} $\rightarrow$ CE};

% =========================== compensation caps =============================
\draw[wire] (\xR,\ycpL) -- (\xR+1.4,\ycpL);
\draw (\xR+1.4,\ycpL) to[C, l=\scriptsize CpL] (\xR+2.6,\ycpL);
\draw[wire] (\xR+2.6,\ycpL) -- (\xB,\ycpL);
\node[nd] at (\xR,\ycpL) {}; \node[nd] at (\xB,\ycpL) {};

\draw[wire] (\xL,\ycpR) -- (\xL+2.9,\ycpR);
\draw (\xL+2.9,\ycpR) to[C, l=\scriptsize CpR] (\xL+4.1,\ycpR);
\draw[wire] (\xL+4.1,\ycpR) -- (\xO,\ycpR);
\node[nd] at (\xL,\ycpR) {}; \node[nd] at (\xO,\ycpR) {};
\end{circuitikz}}
  \caption[{Control amplifier core: a symmetric current-mirror OTA with a PMOS input pair.}]{Control amplifier core: a symmetric current-mirror (mirrored) OTA with a
  PMOS input pair, drawn in the same arrangement as the Virtuoso cellview --- input pair
  and its two unit diode loads on the left, the two output branches on the right. Each
  input device steers part of $I_{\mathrm{tail}}$ into a unit diode ($m$=1); a mirror
  four times larger then restates it, the 	exttt{VG2L} half directly onto 	exttt{VOut}
  and the 	exttt{VG2R} half into node 	exttt{B}, where the PMOS mirror inverts it and
  returns it to 	exttt{VOut} so the two halves add. Both output devices are cascoded.
  	exttt{CpL} and 	exttt{CpR} damp the mirror poles. Identically labelled nets are
  connected; matching dummies are omitted.}
  \label{fig:ctrlamp-core}
\end{figure}
